%% ma: L12-B08-0156
%% lop: 12
%% chuong: 2
%% bai: 8
%% dang: 12.8.51
%% loai: TN4
%% muc: VD
%% dap_an: "B"
%% nguon: "CĐ2.4.pdf giáo viên gửi 27/10/2026 (M07, Lê Minh Khôi – Nguyễn Anh Khôi, Chương II Chủ đề 4 Các bài toán về tâm tỉ cự), Câu 6"
%% kien_thuc: "Bài 8 (biểu thức tọa độ của phép toán vectơ, độ dài vectơ), Bài 6 (điểm thỏa hệ thức vectơ, quy tắc ba điểm), Bài 7 (hình chiếu lên trục tọa độ)"
%% trang_thai: cho-duyet
%% ly_do: "Dạng mới 12.8.51 đề xuất, chờ giáo viên duyệt"
\begin{ex}
Trong không gian $Oxyz$, cho ba điểm $A(-2;1;1)$, $B(4;3;-3)$, $C(5;0;5)$ và $M$ là điểm thuộc trục hoành sao cho $\left|\overrightarrow{MA}+\overrightarrow{MB}+\overrightarrow{MC}\right|$ đạt giá trị nhỏ nhất. Khi đó hoành độ điểm $M$ thuộc khoảng nào sau đây?
\choice{$(-1;1)$}{\True $(1;3)$}{$(3;5)$}{$(5;7)$}
\loigiai{Trọng tâm $G\left(\dfrac73;\dfrac43;1\right)$, $\left|\overrightarrow{MA}+\overrightarrow{MB}+\overrightarrow{MC}\right|=3MG$.

$MG$ nhỏ nhất khi $M$ là hình chiếu của $G$ lên $Ox$: $x_M=\dfrac73\in(1;3)$.}
\end{ex}
